\section{Call the same operation on different word selections}
One column round applies four quarter-rounds to the four columns. Those selections
are disjoint, so changing column zero does not change any input of column one.
The next round uses diagonals that wrap from the right edge to the left edge.
The following table states the selections in \(a,b,c,d\) order; the integers are
word indices, not byte offsets.

\begin{table}[htbp]
  \centering\caption{Word selections in one double-round.}
  \tagpdfsetup{table/header-rows={1}}
  \begin{tabular}{cll}\toprule
    Call & Column round & Diagonal round \\\midrule
    0    & 0,4,8,12     & 0,5,10,15      \\
    1    & 1,5,9,13     & 1,6,11,12      \\
    2    & 2,6,10,14    & 2,7,8,13       \\
    3    & 3,7,11,15    & 3,4,9,14       \\
    \bottomrule\end{tabular}
\end{table}

A column round followed by a diagonal round is a \emph{double-round}. Repeat that
pair ten times. There are twenty rounds and eighty quarter-round calls. Counting
individual calls as complete rounds would execute too little work.

For the scalar implementation, keep the working base in a saved register. Construct
the four argument pointers immediately before each call. The first column at a
working base in \texttt{s1} begins as follows.
\Needspace{11\baselineskip}
\MARSCodeLabel{Select four separated words of column zero.}
\begin{minted}[samepage=true]{mars-idea}
daddiu $a0, $s1, 0
daddiu $a1, $s1, 16
daddiu $a2, $s1, 32
daddiu $a3, $s1, 48
jal chacha_quarter
nop
\end{minted}

The quarter-round can change argument and temporary registers. A saved base and
saved double-round count survive that nested call. Save the block function's incoming
\texttt{ra} before its first \texttt{jal}; otherwise the last nested return address
would replace the address needed to return to the caller.

Write the four column selections first and inspect their output. Then add the four
diagonal selections. For a diagonal pointer, multiply its selected index by four;
for example index ten is byte offset forty. Advancing every argument by four after
a column works across the column group, but cannot express the wrapped diagonals
without additional selection logic. Explicit offsets make this first version easier
to compare with the table.
